% GEP R2 Mock --- LeetCode Sep 2025 on-campus report --- 30 min
\documentclass[11pt,a4paper]{article}
\usepackage[margin=1.55cm,top=1.9cm,bottom=1.9cm]{geometry}
\usepackage[T1]{fontenc}
\usepackage{lmodern}
\usepackage{xcolor}
\usepackage{titlesec}
\usepackage{enumitem}
\usepackage{fancyhdr}
\usepackage{parskip}
\usepackage[hidelinks]{hyperref}
\setlength{\parskip}{1.5pt}

\definecolor{ink}{HTML}{1F2937}
\definecolor{accent}{HTML}{0F5C4C}
\definecolor{muted}{HTML}{6B7280}
\definecolor{soft}{HTML}{F3F6F5}
\definecolor{qcol}{HTML}{9F1239}
\definecolor{you}{HTML}{166534}
\definecolor{youbg}{HTML}{ECFDF5}
\definecolor{rulec}{HTML}{D7E3DE}
\definecolor{tag}{HTML}{B45309}
\definecolor{timec}{HTML}{1D4ED8}

\color{ink}
\pagestyle{fancy}
\fancyhf{}
\fancyhead[L]{\small\color{muted}GEP Intern -- Technology}
\fancyhead[R]{\small\color{accent}\textbf{R2 Mock --- LeetCode}}
\fancyfoot[C]{\small\color{muted}\thepage}
\renewcommand{\headrulewidth}{0.4pt}
\renewcommand{\headrule}{\color{rulec}\hrule height 0.4pt}

\titleformat{\section}{\large\bfseries\color{accent}}{\thesection.}{0.5em}{}
\titlespacing{\section}{0pt}{12pt}{5pt}
\titleformat{\subsection}{\normalsize\bfseries\color{ink}}{}{0em}{}
\titlespacing{\subsection}{0pt}{8pt}{3pt}

\newcommand{\card}[1]{%
  \par\vspace{4pt}%
  \noindent\colorbox{soft}{%
    \parbox{\dimexpr\textwidth-2\fboxsep}{\vspace{5pt}#1\vspace{5pt}}%
  }\par\vspace{5pt}%
}
\newcommand{\IV}[1]{\par\vspace{4pt}\noindent\textcolor{qcol}{\textbf{Interviewer.}}~\textit{#1}\par}
\newcommand{\You}[1]{%
  \par\vspace{2pt}%
  \noindent\colorbox{youbg}{%
    \parbox{\dimexpr\textwidth-2\fboxsep}{\vspace{3pt}\textcolor{you}{\textbf{\small You say.}}~#1\vspace{3pt}}%
  }\par}
\newcommand{\clk}[1]{\noindent\textcolor{timec}{\textbf{#1}}\quad}
\newcommand{\src}[1]{\noindent{\footnotesize\color{tag}\textbf{From report:} #1}\par\vspace{2pt}}
\newcommand{\note}[1]{\noindent{\small\color{muted}\textit{#1}}\par\vspace{2pt}}

\begin{document}

\begin{center}
  {\LARGE\bfseries\color{accent}R2 Mock --- LeetCode 2025 Report}\\[5pt]
  {\color{muted}30 minutes $\cdot$ Technical Interview 2 $\cdot$ Summer Intern 2026}\\[3pt]
  {\footnotesize\color{muted}Mapped to your resume. Verdict in report: Selected, Mumbai, Rs.\ 50k/month.}
\end{center}

\src{LeetCode Discuss, Sep 2025: R2 = brief intro, casual tone, deep resume/projects, design choices + challenges, behavioral (team, conflicts, new stack), balanced technical + cultural. No hard DSA in R2. HR later asked layman project explanation.}

\textit{Interviewer tone: starts casual, puts you at ease, then goes deep on whatever is on the resume.}

\card{
\textbf{How to use.} 30-minute timer. Answer out loud first; green boxes are the target answers mapped to \textbf{your} projects.\\[4pt]
\textbf{R2 intro:} Same 40-second version as Dilip mock (LogiFlow pipelines + Community Hero + Nexus sixteen).\\[4pt]
\textbf{This mock covers R2 only.} R1 topics (SSR, SQL join, max-in-string) appear in the bonus appendix if R1 and R2 blur on your drive day.
}

\section{Clock (matches report structure)}

\begin{tabular}{@{}clp{8.5cm}@{}}
\clk{0:00--3:00} & Intro + rapport & Brief intro; casual chat to settle in. \\
\clk{3:00--12:00} & Resume / LogiFlow & Tech stack, design choices, challenges. \\
\clk{12:00--18:00} & Community Hero & Solo project, AI, auth, trade-offs. \\
\clk{18:00--25:00} & Behavioral & Team conflict, new stack, collaboration. \\
\clk{25:00--30:00} & Wrap-up & Your questions; thank you. \\
\end{tabular}

\section{Script --- Technical Interview 2}

\subsection*{0:00 --- Opening}

\IV{Hi Ojas, how are you doing? Hope round one went okay. I will keep this relaxed --- just want to understand your projects and how you work in a team. Quick intro?}

\You{R2 intro. Stop.}

\IV{Nice. Let us start with LogiFlow. Walk me through the tech stack and \emph{why} you picked each piece.}

\You{Backend Python + FastAPI --- data pipelines and sklearn delay model. Frontend TypeScript + Next.js --- map UI with faster first load. SQL as source of truth for corridors and tariffs. Redis cache-aside for hot reads, 100--400 ms. pytest 100/100 business rules in CI. Deployed on GCP Cloud Run. Each choice matched the constraint: data + ranking on the backend, read-heavy API, tests as contracts.}

\IV{What was the hardest challenge while building it?}

\You{Latency. First version recomputed rankings on every request --- seconds. I moved to cache-aside with Redis and kept SQL as source of truth. That took responses to 100--400 milliseconds. Second challenge was team alignment on what ``co-lead'' owns --- we fixed it by naming pipeline owners in the README.}

\IV{LogiFlow was a team project. Any conflict with a teammate? How did you handle it?}

\You{Cache placement --- API layer vs pipeline layer. We built both on a branch, benchmarked, picked API-level for simpler invalidation with the same latency. No personal conflict; we let data decide.}

\IV{Design choice question: why SQL database, not a document store like MongoDB?}

\You{Corridors, legs, and tariffs are relational --- joins and foreign keys. One tariff update is one row, not thousands of denormalized documents. Mongo fits flexible blobs; my Community Hero reports are exactly that, so that project uses Firestore instead.}

\subsection*{12:00 --- Second project (they often switch)}

\IV{Community Hero --- you said you built it solo. Tell me about that. Tech stack and why.}

\You{Civic issue reporting PWA. React + Express, Firestore for geo-tagged reports, Cloud Run deploy, Gemini for category classification. Solo because hackathon timeline --- I owned spec to deploy. Firestore because each report is a self-contained document with nested fields; schema was still moving. Gemini behind guard checks --- bad output falls back to manual admin review.}

\IV{What if Gemini is wrong or the API is down?}

\You{Treat the model as untrusted. Validate input before the call. If output is not a known category, store raw text and flag for admin. Prefer a missed auto-tag over a wrong auto-route. If Gemini is down, reports still save; routing goes to manual queue.}

\IV{How did you handle login and authentication?}

\You{Firebase Auth on the client issues a token. Express verifies it on every protected route before Firestore. Roles checked server-side --- a citizen cannot hit admin endpoints. I never trust the client for authorisation.}

\subsection*{18:00 --- Behavioral (quoted in report)}

\IV{Would you prefer to keep working in Python and React, or are you open to a completely new tech stack in a corporate setting?}

\You{Open to a new stack. GEP often uses C\#, Angular, SQL Server --- I have not shipped those, but I learned FastAPI, Redis, and Cloud Run under deadline for LogiFlow. The habits transfer: APIs, data, tests, ownership. I would rather learn the product stack than insist on mine.}

\IV{How do you collaborate on tasks when the team is under deadline?}

\You{Name owners early --- on LogiFlow we split pipelines by person and shared a ranker interface. Daily sync on blockers, PR reviews before merge. On Nexus, sixteen-person core team --- divide events, one owner per deliverable. I do not merge what I have not reviewed.}

\IV{If I asked you to explain one project to someone with zero technical background, which would you pick and how?}

\You{LogiFlow. ``You need to ship goods from city A to B. Train, truck, or both --- different price, time, and delay risk. LogiFlow shows all options ranked so you pick what fits your budget and deadline --- like Google Maps for freight, but you see trade-offs, not just one fastest route.''}

\note{Report says this layman question appeared in \textbf{HR}; it may come in R2 if time allows. Have it ready either way.}

\subsection*{25:00 --- Wrap-up}

\IV{That is most of what I had. Anything you want to ask me?}

\You{For technology interns, what does the first month look like --- shadowing, or contributing to a live module? And what separates interns who get a PPO from those who do not?}

\IV{Good questions. Thank you, Ojas --- HR will connect with you on next steps.}

\You{Thank you. I enjoyed this.}

\section{Bonus --- if R1 topics repeat (from same LeetCode report)}

\card{\textbf{R1 was 30 min in the report.} If your drive merges rounds or Dilip asks R1-style questions, drill these once:}

\IV{What is SSR in Next.js?}
\You{Server renders React to HTML before sending to the browser, so the user sees content immediately. Plain React sends empty HTML + JS first, then fetches data --- blank screen longer. LogiFlow map benefits from data in the first paint.}

\IV{Stack vs Queue?}
\You{Stack LIFO --- undo, DFS, call stack. Queue FIFO --- BFS, job scheduling. LogiFlow graph search uses a queue for BFS on corridors.}

\IV{SQL: books and authors --- find books by author X.}
\You{\texttt{SELECT b.name FROM book b JOIN author a ON b.author\_id = a.author\_id WHERE a.name = 'X';}}

\IV{String s = ``1,2,3,4,5'' --- find maximum integer.}
\You{Split on comma, parse each token to int, track max in one pass. O(n) time, O(1) extra if split in place. Watch empty string and non-numeric tokens --- ask if input is guaranteed clean.}

\section{Score this mock}

\begin{tabular}{@{}p{3.2cm}p{9cm}c@{}}
\textbf{Rapport} & Short intro, not a monologue? Comfortable tone? & \rule{1.4cm}{0.4pt} \\
\textbf{LogiFlow} & Stack + why each? Challenge + fix? SQL vs NoSQL? & \rule{1.4cm}{0.4pt} \\
\textbf{Community Hero} & Solo ownership clear? Gemini guards + fallback? & \rule{1.4cm}{0.4pt} \\
\textbf{Behavioral} & New stack = yes + example? Conflict = evidence? & \rule{1.4cm}{0.4pt} \\
\textbf{Layman} & No jargon in dinner-table explanation? & \rule{1.4cm}{0.4pt} \\
\textbf{Close} & Two thoughtful questions, not stipend? & \rule{1.4cm}{0.4pt} \\
\end{tabular}

\vspace{8pt}
\card{\textbf{Report reminder.} Selected candidate: R2 had \textbf{no} heavy coding. Real filter was project depth + team fit. HR asked family + layman project --- prep those even if R2 feels ``easy.''\\[4pt]
Pair with: \texttt{GEP\_R2\_Mock\_Dilip.pdf} and \texttt{GEP\_R2\_LogiFlow\_Prep.pdf}.}

\end{document}
